\documentclass[11pt]{amsart}
\usepackage[T1]{fontenc}
\usepackage{lmodern}
\usepackage[margin=1.08in]{geometry}
\usepackage{microtype}
\usepackage{amssymb}
\usepackage{mathtools}
\usepackage[colorlinks=true,linkcolor=blue,citecolor=blue,urlcolor=blue]{hyperref}
\hypersetup{pdftitle={Almost everywhere monotone additive functions},pdfauthor={Qiyuan Gu}}
\numberwithin{equation}{section}
\newtheorem{theorem}{Theorem}[section]
\newtheorem{lemma}[theorem]{Lemma}
\newtheorem{proposition}[theorem]{Proposition}
\theoremstyle{remark}
\newtheorem{remark}[theorem]{Remark}
\newcommand{\N}{\mathbb N}
\newcommand{\R}{\mathbb R}
\newcommand{\E}{\mathbb E}
\newcommand{\ind}{\mathbf 1}
\newcommand{\clip}{\phi_K}

\title[Monotonicity of additive functions]{Almost everywhere monotone additive functions}
\author{Qiyuan Gu}
\address{University of Chicago}
\email{phoenix1203@uchicago.edu}
\date{September 7, 2026}

\begin{document}
\begin{abstract}
We prove that a real additive function which decreases on a set of
natural density zero is a nonnegative constant multiple of the logarithm.
This answers a conjecture of Erd\H{o}s from 1946, listed as Problem~1122
on the Erd\H{o}s Problems website. The main step is to deduce finite
concentration without a growth assumption on the values at prime powers.
A truncated normalization and a bounded clipping reduce this step to
Mangerel's first-moment theorem for short intervals and the moment
estimates of Ruzsa and Hildebrand.
\end{abstract}
\maketitle

\section{Introduction}

A function $f\colon\N\to\R$ is \emph{additive} if
$f(ab)=f(a)+f(b)$ whenever $(a,b)=1$. Erd\H{o}s proved that an additive
function which is nondecreasing on $\N$ must be a constant multiple of
$\log n$. He asked whether the same conclusion follows when the set of
decreases has natural density zero \cite[p.~3]{Erdos}. The question is
also recorded as Erd\H{o}s Problem~1122 \cite{Problem1122}.

\begin{theorem}\label{thm:main}
Let $f\colon\N\to\R$ be additive. If
\[
 D_f(X):=\#\{n\le X:f(n+1)<f(n)\}=o(X),
\]
then there is a constant $c\ge0$ such that $f(n)=c\log n$ for every
$n\in\N$.
\end{theorem}

Mangerel established an approximate logarithmic representation under the
same density hypothesis \cite[Theorem~1.9]{Mangerel}. For completely
additive functions, he obtained the exact conclusion with additional
assumptions on the values at primes and the number of decreases
\cite[Corollary~1.7]{Mangerel}. His short-interval theorem is the
analytic input used here.

We say that $f$ has \emph{finite concentration} if there are constants
$d>0$ and $R<\infty$ such that, for arbitrarily large $X$, some interval
of length $R$ contains at least $dX$ of the values $f(n)$, $n\le X$.
The interval is allowed to depend on $X$. Two results of Erd\H{o}s reduce
Theorem~\ref{thm:main} to finite concentration. Theorem~V of
\cite{Erdos} states that finite concentration is equivalent to the
existence of $c\in\R$ for which
\begin{equation}\label{eq:classical-concentration}
 \sum_p\frac{\min\{(f(p)-c\log p)^2,1\}}p<\infty.
\end{equation}
Theorem~X of the same paper implies that, under
\eqref{eq:classical-concentration}, the limiting distribution of
$f(n+1)-f(n)$ has no mass on $(-\infty,0)$ only when
$f(n)=c\log n$ for all $n$. Both results apply to general additive
functions, with no restriction on their values at higher prime powers.

The proof below establishes finite concentration by contradiction.
If it fails, we normalize $f$ and subtract a logarithm so that the
truncated quadratic mass of the prime values is a small fixed number
$M$. The minimizing logarithmic coefficient supplies an orthogonality
identity. We then compare a bounded clipping of the normalized function
with a strongly additive function obtained by clipping its prime values.
Their mean-square difference is small enough to retain a positive
variance. On the other hand, the density hypothesis makes the bounded
function almost constant on a typical interval of any fixed length.
Short-interval averaging forces the variance of the strongly additive
function to be smaller than its moment lower bound permits.

\subsection*{Notation}
Letters $p,q$ denote primes. We write
\[
 \E_X F=\frac1X\sum_{n\le X}F(n),\qquad
 \mathcal A_H F(n)=\frac1H\sum_{j=0}^{H-1}F(n-j)\quad(n\ge H).
\]
Restricted averages, such as $\E_X(F\ind_{n\ge H})$, retain the
normalizing factor $1/X$. For a strongly additive function
$z(n)=\sum_{p\mid n}a_p$, put
\[
 \mu_z(X)=\sum_{p\le X}\frac{a_p}{p}.
\]
All implicit constants are absolute unless a dependence is indicated.
The constants $K$ and $M$ used in the proof are fixed before $X$ tends
to infinity. In statements about a family $z_X$, the prime coefficients
may depend on $X$.

\section{Moment estimates and short intervals}

For an additive function $z$, define
\[
 A_0(z,X)=\sum_{p^k\le X}\frac{z(p^k)}{p^k},\qquad
 B_j(z,X)^j=\sum_{p^k\le X}\frac{|z(p^k)|^j}{p^k}\quad(j=2,4).
\]
We use the following consequences of Hildebrand's moment theorem
\cite[Th\'eor\`eme~$1^*$, equation~(6)]{Hildebrand}:
\begin{align}
 \E_X|z-A_0(z,X)|^2
 &\asymp\inf_{b\in\R}\left(b^2+
       \sum_{p^k\le X}\frac{|z(p^k)-b\log(p^k)|^2}{p^k}\right),
       \label{eq:ruzsa}\\
 \E_X|z-A_0(z,X)|^4
 &\ll B_2(z,X)^4+B_4(z,X)^4.\label{eq:fourth}
\end{align}
The second-moment equivalence \eqref{eq:ruzsa} is due to Ruzsa;
Hildebrand's theorem includes it as the case of exponent~$2$.

If $z$ is strongly additive, then
\begin{equation}\label{eq:center-difference}
 |A_0(z,X)-\mu_z(X)|
 \le\sum_{p\le X}\frac{|a_p|}{p(p-1)}
 \ll\left(\sum_{p\le X}\frac{a_p^2}p\right)^{1/2}.
\end{equation}
Also $B_2(z,X)^2\le2\sum_{p\le X}a_p^2/p$, with the analogous
bound for $B_4$. Thus, if $|a_p|\le L$ and
$\sum_{p\le X}a_p^2/p\le V$, then
\begin{equation}\label{eq:bounded-moments}
 \E_X|z-\mu_z(X)|^2\ll V,\qquad
 \E_X|z-\mu_z(X)|^4\ll V^2+L^2V.
\end{equation}
If the coefficients are bounded and $a_p\to0$ for each fixed prime,
dominated convergence in \eqref{eq:center-difference} gives
$A_0(z,X)-\mu_z(X)\to0$.

\begin{lemma}\label{lem:short}
Let $z_X$ be strongly additive functions satisfying
\[
 |z_X(p)|\le L,\qquad
 \sum_{p\le X}\frac{z_X(p)^2}{p}\le V
\]
for fixed $L,V<\infty$. Then
\begin{equation}\label{eq:short-limit}
 \lim_{H\to\infty}\limsup_{X\to\infty}
 \frac1X\sum_{H\le n\le X}
 |\mathcal A_Hz_X(n)-\mu_{z_X}(X)|^2=0,
\end{equation}
where $H$ runs through the positive integers.
\end{lemma}

\begin{proof}
For $Y\le X$, write
\[
 d_Y=\frac2Y\sum_{Y/2<n\le Y}z_X(n),\qquad
 F_Y(n)=\mathcal A_Hz_X(n)-d_Y.
\]
Mangerel's Theorem~1.1 \cite{Mangerel} gives, for every integer
$10\le H\le Y/100$,
\begin{equation}\label{eq:mangerel}
 \frac2Y\sum_{Y/2<n\le Y}|F_Y(n)|
 \ll\left(\sqrt{\frac{\log\log H}{\log H}}
                +(\log Y)^{-1/800}\right)\sqrt{2V}.
\end{equation}
The estimate is uniform over all additive functions; in particular, it
applies to the changing family $z_X$.

Counting multiples of each prime gives
\begin{equation}\label{eq:dyadic-center}
 d_Y-\mu_{z_X}(Y)
 =\sum_{p\le Y}z_X(p)
 \left(\frac2Y\left(\left\lfloor\frac Yp\right\rfloor
       -\left\lfloor\frac Y{2p}\right\rfloor\right)-\frac1p\right)
 =O\left(\frac{L\pi(Y)}Y\right).
\end{equation}
Jensen's inequality and \eqref{eq:bounded-moments} at $Y$ imply
\begin{equation}\label{eq:discrepancy-fourth}
 \frac2Y\sum_{Y/2<n\le Y}|F_Y(n)|^4
 \ll V^2+L^2V+O_L((\log Y)^{-4}).
\end{equation}
To see this, center the backward average at $\mu_{z_X}(Y)$.
Every index in its inner sums lies in $[1,Y]$ and occurs at most $H$
times. The remaining constant is estimated by
\eqref{eq:dyadic-center}.
H\"older's inequality now gives
\[
 \frac2Y\sum|F_Y|^2
 \le\left(\frac2Y\sum|F_Y|\right)^{2/3}
       \left(\frac2Y\sum|F_Y|^4\right)^{1/3}.
\]
Together with \eqref{eq:mangerel} and
\eqref{eq:discrepancy-fourth}, this proves the required iterated limit
on each fixed dyadic band $Y\asymp X$.

For completeness, this dyadic conclusion extends to the whole interval
as follows. On each of finitely many bands with $Y\asymp X$, we have
\[
 |\mu_{z_X}(X)-\mu_{z_X}(Y)|
 \le L\sum_{Y<p\le X}\frac1p=o(1).
\]
The omitted initial segment $H\le n\le\eta X$ contributes
$O_{L,V}(\sqrt\eta)$ to \eqref{eq:short-limit}: apply
Cauchy--Schwarz and Jensen, using the fourth-moment bound
\eqref{eq:bounded-moments} at $X$. First take $X\to\infty$,
then $H\to\infty$, and finally $\eta\downarrow0$.
\end{proof}

\section{A weighted estimate for a set of primes}

\begin{lemma}\label{lem:tail-weight}
Suppose $0<M<1/4$ and $\mathcal T$ is a set of primes at most $X$
such that $\sum_{p\in\mathcal T}1/p\le M$. Then, uniformly in
$\mathcal T$,
\begin{equation}\label{eq:tail-weight}
 \sum_{p\in\mathcal T}\frac1p
       \sum_{X/p<q\le X}\frac1q
 \ll M\log(2/M)+o_{X\to\infty}(1).
\end{equation}
Here $M$ is fixed in the term $o(1)$.
\end{lemma}

\begin{proof}
For $p\le X^{1-M}$, the inner sum is at most
$\log(1/M)+o(1)$ by the prime harmonic-sum estimate. These primes
contribute at most $M\log(1/M)+o(1)$.

For $p>X^{1-M}$, discard the restriction $p\in\mathcal T$ and
interchange the sums. The terms with $q>X^M$ are bounded by
\[
 \left(\sum_{X^{1-M}<p\le X}\frac1p\right)
 \left(\sum_{X^M<q\le X}\frac1q\right)
 \ll M\log(2/M)+o(1).
\]
For $q\le X^M$, the condition $X/p<q$ implies $p>X/q$.
The bound $\pi(t)\ll t/\log t$ and partial summation give,
uniformly for $2\le q\le X^M$,
\[
 \sum_{X/q<p\le X}\frac1p\ll\frac{\log q}{\log X}.
\]
Their total contribution is therefore
\[
 \ll\frac1{\log X}\sum_{q\le X^M}\frac{\log q}{q}\ll M.
\]
This proves \eqref{eq:tail-weight}.
\end{proof}

\section{Normalization}

In the remainder of the proof, suppose that $f$ does not have finite
concentration. For $s>0$ define
\begin{equation}\label{eq:W}
 W_X(s)=\min_{c\in\R}
 \left\{c^2+\sum_{p\le X}
       \frac{\min\{(f(p)/s-c\log p)^2,1\}}p\right\}.
\end{equation}

\begin{lemma}\label{lem:normalization}
For each fixed $M>0$ and all sufficiently large $X$, there are
$s_X>0$ and a minimizer $c_X$ in \eqref{eq:W} such that
$W_X(s_X)=M$. They may be chosen so that $s_X\to\infty$; any
such choices satisfy $c_X\to0$. If
\[
 u_p=\frac{f(p)}{s_X}-c_X\log p\quad(p\le X),
\]
then
\begin{align}
 \sum_{p\le X}\frac{\min\{u_p^2,1\}}p
 &=M-c_X^2=M-o(1),\label{eq:mass}\\
 \sum_{\substack{p\le X\\|u_p|<1}}
       \frac{u_p\log p}{p}&=c_X.\label{eq:stationarity}
\end{align}
No $u_p$ has absolute value exactly $1$.
\end{lemma}

\begin{proof}
The expression in braces in \eqref{eq:W} is continuous and coercive
as a function of $c$, so its minimum is attained. The minimum is
continuous in $s>0$, since minimizers remain bounded when $s$ varies
over a compact subinterval. On writing $\lambda=sc$, the expression
becomes
\[
 \frac{\lambda^2}{s^2}
 +\sum_{p\le X}\frac{\min\{(f(p)-\lambda\log p)^2/s^2,1\}}p.
\]
It follows that $W_X(s)$ is nonincreasing in $s$, and it tends to
zero as $s\to\infty$ for fixed $X$.

For each fixed $s>0$, we claim that $W_X(s)\to\infty$. Otherwise
some unbounded sequence of $X$ would have bounded $W_X(s)$ and
bounded minimizers. Pass to a subsequence along which the minimizers
converge to $c$. Fatou's lemma gives
\[
 \sum_p\frac{\min\{(f(p)/s-c\log p)^2,1\}}p<\infty.
\]
Rescaling by the fixed number $s$ preserves convergence of the
truncated square sum. Erd\H{o}s's Theorem~V would therefore give
finite concentration of $f$, contrary to the assumption.

Continuity now supplies $s_X$ with $W_X(s_X)=M$. Monotonicity and
the preceding claim imply $s_X\to\infty$. We have
$|c_X|\le\sqrt M$. If a subsequential limit of $c_X$ were a
nonzero number $c$, Fatou's lemma on the fixed primes would give
\[
 \sum_p\frac{\min\{(c\log p)^2,1\}}p\le M,
\]
which is impossible. Thus $c_X\to0$, and \eqref{eq:mass} follows.

It remains to justify differentiation at the minimizer. At a point
where $|f(p)/s_X-c\log p|=1$, the derivative of the corresponding
summand has a downward jump of size $2\log p/p$ as $c$ increases.
Coincident jumps add. A local minimum of a continuous piecewise
differentiable function has left derivative at most zero and right
derivative at least zero, so a downward jump excludes a local minimum.
Consequently no such point is a minimizer. Differentiating
\eqref{eq:W} at $c_X$ gives \eqref{eq:stationarity}.
\end{proof}

\section{Clipping the normalized function}

Fix $K>1$ and $0<M<1/4$, and use Lemma~\ref{lem:normalization}.
Let $\clip(t)=\max\{-K,\min\{t,K\}\}$ and set
\begin{align*}
 \mathcal T&=\{p\le X:|u_p|>1\},&
 S_0(n)&=\sum_{\substack{p\mid n\\p\notin\mathcal T}}u_p,\\
 a_X&=\sum_{\substack{p\le X\\p\notin\mathcal T}}\frac{u_p}{p},&
 S(n)&=S_0(n)-a_X.
\end{align*}
Coefficients at primes greater than $X$ are set to zero. Define
\begin{equation}\label{eq:Y-and-Z}
 \begin{split}
 z_X(n)&=\sum_{p\mid n}\clip(u_p),\qquad Z(n)=z_X(n)-a_X,\\
 Y(n)&=\clip\left(\frac{f(n)}{s_X}-c_X\log n-a_X\right).
 \end{split}
\end{equation}
The distinction is that $Y$ clips the value of the additive function,
whereas $z_X$ clips each prime value before adding.

\begin{lemma}\label{lem:clipping}
There is an absolute constant $C$ such that
\begin{equation}\label{eq:clipping-bound}
 \limsup_{X\to\infty}\E_X|Y-Z|^2
 \le C\left(\frac M{K^2}+M^2\log(2/M)+K^2M^2\right).
\end{equation}
\end{lemma}

\begin{proof}
The small coefficients have quadratic mass at most $M$, and
$\sum_{p\in\mathcal T}1/p\le M$. Hence
\begin{equation}\label{eq:S-fourth}
 \E_X S^4\ll M+M^2.
\end{equation}
Put $T(n)=\#\{p\in\mathcal T:p\mid n\}$. Counting multiples of
distinct primes gives
\begin{equation}\label{eq:tail-count}
 \E_X T(T-1)
 =\sum_{\substack{p,q\in\mathcal T\\p\ne q}}
       \frac1X\left\lfloor\frac X{pq}\right\rfloor
 \le M^2.
\end{equation}
We also need
\begin{equation}\label{eq:mixed-moment}
 \E_X\bigl(S^2\ind_{T\ge1}\bigr)
 \ll M^2\log(2/M)+o(1).
\end{equation}
For $p\in\mathcal T$, the coefficient of $p$ in $S_0$ is zero,
so $S_0(pm)=S_0(m)$ for every $m$, including multiples of $p$.
Writing $y=X/p$ and using \eqref{eq:bounded-moments} at $y$, we get
\begin{align*}
 \E_X\bigl(S^2\ind_{p\mid n}\bigr)
 &=\frac1p\frac1y\sum_{m\le y}|S_0(m)-a_X|^2\\
 &\ll\frac1p\left(M+
           |a_X-\mu_{S_0}(y)|^2\right).
\end{align*}
At $y=1$, we have $S_0(1)=\mu_{S_0}(1)=0$, so the same bound follows
from the mean-shift term. By Cauchy--Schwarz,
\[
 |a_X-\mu_{S_0}(X/p)|^2
 \le M\sum_{X/p<q\le X}\frac1q.
\]
Now use $\ind_{T\ge1}\le\sum_{p\in\mathcal T}\ind_{p\mid n}$,
sum over $p\in\mathcal T$, and apply
Lemma~\ref{lem:tail-weight}. This proves \eqref{eq:mixed-moment}.

First replace the input to the clipping by the sum of its prime
contributions, and write
\[
 Y_*(n)=\clip\left(S(n)+
                  \sum_{\substack{p\mid n\\p\in\mathcal T}}u_p\right).
\]
If $T(n)=0$, the difference $Y_*(n)-Z(n)$ vanishes for $|S(n)|\le K$
and otherwise has absolute value at most $|S(n)|$. Its squared
contribution is at most $\E_X S^4/K^2$.
If $T(n)=1$, with unique tail value $v$, the Lipschitz property of
$\clip$ gives
\[
 |\clip(S+v)-S-\clip(v)|\le2|S|.
\]
If $T(n)\ge2$, boundedness of $\clip$ gives
\[
 |Y_*(n)-Z(n)|^2\le C S(n)^2+C K^2T(n)(T(n)-1).
\]
Equations \eqref{eq:S-fourth}--\eqref{eq:mixed-moment} therefore imply
\begin{equation}\label{eq:strong-clipping}
 \E_X|Y_*-Z|^2
 \ll\frac M{K^2}+M^2\log(2/M)+K^2M^2+o(1).
\end{equation}

We next compare $Y$ and $Y_*$. If $p^k\Vert n$ with $k\ge2$, its
contribution to the difference between the two unclipped expressions is
\[
 \frac{f(p^k)-f(p)}{s_X}-c_X(k-1)\log p.
\]
This tends to zero for each fixed prime power. Moreover,
$\sum_{p,k\ge2}p^{-k}<\infty$. Given $\varepsilon>0$, choose a
finite set of prime powers so that the sum of $p^{-k}$ outside the set
is less than $\varepsilon$. Uniformly in $X$, the proportion of
integers divisible by any of those excluded prime powers is at most
$\varepsilon$. On the complement, only the fixed finite set can
contribute, and its total contribution tends uniformly to zero.
Thus the unclipped difference tends to zero in probability.
The function $\clip$ is bounded and Lipschitz, so
\[
 \E_X|Y-Y_*|^2=o(1).
\]
Together with \eqref{eq:strong-clipping}, this proves the lemma.
\end{proof}

\begin{lemma}\label{lem:variance-lower}
There are absolute constants $c_0,C_0>0$ such that
\begin{equation}\label{eq:variance-lower}
 \liminf_{X\to\infty}\E_X|z_X-\mu_{z_X}(X)|^2
 \ge c_0M-C_0K^2M^2.
\end{equation}
\end{lemma}

\begin{proof}
Write $v_p=\clip(u_p)$. Since $K>1$, \eqref{eq:mass} implies
\begin{equation}\label{eq:clipped-mass}
 M-o(1)\le\sum_{p\le X}\frac{v_p^2}{p}\le K^2M.
\end{equation}
On $p\notin\mathcal T$ we have $v_p=u_p$, so
\eqref{eq:stationarity} gives
\begin{equation}\label{eq:projection-numerator}
 \left|\sum_{p\le X}\frac{v_p\log p}{p}\right|
 \le |c_X|+KM\log X.
\end{equation}
Using $\sum_{p\le X}(\log p)^2/p\sim\tfrac12(\log X)^2$, we obtain
\begin{align}
 \inf_{b\in\R}\sum_{p\le X}\frac{(v_p-b\log p)^2}{p}
 &=\sum_{p\le X}\frac{v_p^2}{p}
   -\frac{\bigl(\sum_{p\le X}v_p\log p/p\bigr)^2}
          {\sum_{p\le X}(\log p)^2/p}\notag\\
 &\ge M-C K^2M^2-o(1).\label{eq:projection}
\end{align}
Apply \eqref{eq:ruzsa}, dropping the nonnegative $b^2$ term and the
terms with $k\ge2$ from its right-hand side.
For each fixed prime $p$, $v_p\to0$, because $s_X\to\infty$ and
$c_X\to0$. Therefore \eqref{eq:center-difference} and dominated
convergence give $A_0(z_X,X)-\mu_{z_X}(X)\to0$.
The second moments are uniformly bounded by
\eqref{eq:bounded-moments} and \eqref{eq:clipped-mass}; changing the
center consequently changes the second moment by $o(1)$.
Combining this with \eqref{eq:projection} proves
\eqref{eq:variance-lower}.
\end{proof}

\section{Proof of Theorem~\ref{thm:main}}

Suppose first, as in Sections~4 and~5, that $f$ has no finite
concentration. The function $Y$ in \eqref{eq:Y-and-Z} is bounded
between $-K$ and $K$. If $f(n+1)\ge f(n)$, monotonicity and
Lipschitz continuity of $\clip$ imply
\[
 (Y(n)-Y(n+1))_+\le |c_X|\log(1+1/n),
\]
where $t_+=\max\{t,0\}$. At a decrease of $f$, the same quantity
is at most $2K$. Telescoping the signed increments gives
\begin{equation}\label{eq:variation}
 \sum_{n<X}|Y(n+1)-Y(n)|
 \le2K+4KD_f(X)+2|c_X|\log X=o(X).
\end{equation}
For each fixed $H$, every difference $Y(n)-Y(n-j)$ is bounded by
the sum of the $j$ intervening absolute increments. Since $|Y|\le K$,
\eqref{eq:variation} yields
\begin{equation}\label{eq:Y-short}
 \frac1X\sum_{H\le n\le X}|Y(n)-\mathcal A_HY(n)|^2\longrightarrow0.
\end{equation}

Let $m_X=\mu_{z_X}(X)-a_X$. Then $Z-m_X=z_X-\mu_{z_X}(X)$, and
\begin{equation}\label{eq:four-terms}
 Z-m_X=(Z-Y)+(Y-\mathcal A_HY)
            +\mathcal A_H(Y-Z)+(\mathcal A_H Z-m_X).
\end{equation}
The square of the right-hand side is at most four times the sum of
the four squared terms. Jensen's inequality gives the contraction
\[
 \frac1X\sum_{H\le n\le X}|\mathcal A_H(Y-Z)(n)|^2
 \le\frac1X\sum_{n\le X}|Y(n)-Z(n)|^2.
\]
The prime coefficients of $z_X$ are bounded by $K$ and have quadratic
mass at most $K^2M$, so Lemma~\ref{lem:short} applies to the last
term of \eqref{eq:four-terms}. The first $H$ indices contribute $o(1)$
to the centered second moment: this follows directly from
\eqref{eq:bounded-moments} and Cauchy--Schwarz. Taking $X\to\infty$
and then $H\to\infty$ in \eqref{eq:four-terms}, and using
\eqref{eq:Y-short}, we obtain
\begin{align}
 \limsup_X\E_X|z_X-\mu_{z_X}(X)|^2
 &\le8\limsup_X\E_X|Y-Z|^2\notag\\
 &\le C_1\left(\frac M{K^2}+M^2\log(2/M)+K^2M^2\right)
       \label{eq:variance-upper}
\end{align}
for an absolute constant $C_1$.

Choose $K>1$ so large that $C_1/K^2<c_0/4$, where $c_0$ is from
Lemma~\ref{lem:variance-lower}. Next choose $0<M<1/4$ so small that
\[
 C_1M\log(2/M)+(C_1+C_0)K^2M<c_0/4.
\]
With these fixed choices, \eqref{eq:variance-upper} contradicts
\eqref{eq:variance-lower}. Hence $f$ has finite concentration.

Erd\H{o}s's Theorem~V now gives a constant $c$ for which the additive
remainder $r(n)=f(n)-c\log n$ satisfies
\[
 \sum_p\frac{\min\{r(p)^2,1\}}p<\infty.
\]
By his Theorem~X, the differences $f(n+1)-f(n)$ have a limiting
distribution, and this distribution has no negative mass only if
$r$ is identically zero. For each $\varepsilon>0$, our hypothesis
gives
\[
 \frac1X\#\{n\le X:f(n+1)-f(n)<-\varepsilon\}
 \le\frac{D_f(X)}X\longrightarrow0.
\]
The limiting distribution therefore has no mass on $(-\infty,0)$.
It follows that $f(n)=c\log n$ for all $n$. Finally, $c<0$ would
give a decrease at every integer, so $c\ge0$.\qed

\section*{Declaration of generative AI use}
OpenAI GPT-6 Astra, accessed through Codex, was used to develop the
mathematical arguments and draft this manuscript. A separate instance
of the same model reviewed the proof and its use of the cited results.
This draft is provided for human review; the model review is not a
substitute for review by the author or for peer review.

\begin{thebibliography}{9}

\bibitem{Erdos}
P.~Erd\H{o}s,
\emph{On the distribution function of additive functions},
Ann. of Math. (2) \textbf{47} (1946), no.~1, 1--20.
\url{https://combinatorica.hu/~p_erdos/1946-06.pdf}.

\bibitem{Hildebrand}
A.~Hildebrand,
\emph{Sur les moments d'une fonction additive},
Ann. Inst. Fourier (Grenoble) \textbf{33} (1983), no.~3, 1--22.
\url{https://www.numdam.org/item/AIF_1983__33_3_1_0/}.

\bibitem{Mangerel}
A.~P.~Mangerel,
\emph{Additive functions in short intervals, gaps and a conjecture of Erd\H{o}s},
Ramanujan J. \textbf{59} (2022), 1023--1090.
\href{https://doi.org/10.1007/s11139-022-00623-y}{doi:10.1007/s11139-022-00623-y}.
Preprint: \href{https://arxiv.org/abs/2108.12351}{arXiv:2108.12351}.

\bibitem{Problem1122}
\emph{Erd\H{o}s Problem \#1122}, Erd\H{o}s Problems website,
\url{https://www.erdosproblems.com/1122}, accessed September~7, 2026.

\end{thebibliography}
\end{document}
